\chapter{La parábola}

\begin{center}\begin{minipage}{.88\textwidth}\small\sffamily
\textbf{Propósito.} Interpretar la parábola como lugar geométrico y obtener
sus elementos y ecuaciones en orientación horizontal o vertical.
\end{minipage}\end{center}

\begin{definition}{Parábola}{parabola}
Es el lugar geométrico de los puntos que están a la misma distancia de un
punto fijo, llamado foco, y de una recta fija, llamada directriz.
\end{definition}

El punto medio entre el foco y la directriz sobre el eje de simetría es el
vértice. La distancia dirigida del vértice al foco se denota por $p$.

\begin{formula}[title={Parábolas con vértice en el origen}]
\[
y^2=4px\quad\text{(eje horizontal)},
\qquad
x^2=4py\quad\text{(eje vertical)}.
\]
\end{formula}

\begin{center}
\begin{tabular}{ccll}
\toprule
Ecuación & Foco & Directriz & Apertura\\
\midrule
$y^2=4px$ & $(p,0)$ & $x=-p$ & derecha si $p>0$\\
$x^2=4py$ & $(0,p)$ & $y=-p$ & arriba si $p>0$\\
\bottomrule
\end{tabular}
\end{center}
Si $p<0$, la apertura se invierte.

\section{Deducción de la ecuación}

Para foco $F(p,0)$ y directriz $x=-p$, un punto $P(x,y)$ cumple
\[
PF=\operatorname{dist}(P,\text{directriz}).
\]
Así,
\[
\sqrt{(x-p)^2+y^2}=|x+p|.
\]
Al elevar al cuadrado y simplificar resulta $y^2=4px$.

\begin{example}{Elementos desde la ecuación}{elementos-parabola}
Determina los elementos de $y^2=12x$.
\begin{solution}
\textbf{Paso 1.} Comparamos $4p=12$, por lo que $p=3$.

\textbf{Paso 2.} El vértice es $(0,0)$ y el foco es $(3,0)$.

\textbf{Paso 3.} La directriz es $x=-3$ y el eje es $y=0$.

\textbf{Paso 4.} Como $p>0$, abre hacia la derecha. El lado recto mide
$|4p|=12$.
\end{solution}
\end{example}

\section{Parábolas trasladadas}

\begin{formula}[title={Formas ordinarias}]
\[
(y-k)^2=4p(x-h),
\qquad
(x-h)^2=4p(y-k).
\]
El vértice es $V(h,k)$.
\end{formula}

Para $(y-k)^2=4p(x-h)$:
\[
F(h+p,k),\qquad x=h-p,\qquad y=k.
\]
Para $(x-h)^2=4p(y-k)$:
\[
F(h,k+p),\qquad y=k-p,\qquad x=h.
\]

\begin{example}{Parábola horizontal trasladada}{parabola-horizontal}
Analiza $(y+2)^2=-8(x-1)$.
\begin{solution}
\textbf{Paso 1.} $h=1$, $k=-2$ y $4p=-8$, de modo que $p=-2$.

\textbf{Paso 2.} $V(1,-2)$ y $F(h+p,k)=(-1,-2)$.

\textbf{Paso 3.} La directriz es $x=h-p=3$ y el eje es $y=-2$.

\textbf{Paso 4.} Abre hacia la izquierda; su lado recto mide $8$.
\end{solution}
\end{example}

\begin{example}{Ecuación a partir de foco y vértice}{ecuacion-parabola}
Obtén la parábola con vértice $V(2,1)$ y foco $F(2,5)$.
\begin{solution}
\textbf{Paso 1.} El foco comparte la coordenada $x$ con el vértice, así que
el eje es vertical.

\textbf{Paso 2.} $p=5-1=4$; abre hacia arriba.

\textbf{Paso 3.} Sustituimos en $(x-h)^2=4p(y-k)$:
\[
(x-2)^2=16(y-1).
\]
La directriz es $y=1-4=-3$.
\end{solution}
\end{example}

\section{Lado recto}

El lado recto es la cuerda perpendicular al eje que pasa por el foco. Su
longitud es $|4p|$. Para una parábola horizontal sus extremos son
\[
(h+p,k+2p),\qquad(h+p,k-2p),
\]
y para una vertical son
\[
(h+2p,k+p),\qquad(h-2p,k+p).
\]

\section{De la forma general a la ordinaria}

Una parábola no rotada tiene solamente uno de los términos cuadráticos. Para
obtener su forma ordinaria se completa el cuadrado en la variable elevada al
cuadrado.

\begin{example}{Completar cuadrados}{general-parabola}
Analiza $x^2-6x-8y+1=0$.
\begin{solution}
\textbf{Paso 1.} Agrupamos y trasladamos términos:
\[
x^2-6x=8y-1.
\]
\textbf{Paso 2.} Sumamos $9$ a ambos lados:
\[
(x-3)^2=8y+8=8(y+1).
\]
\textbf{Paso 3.} $h=3$, $k=-1$ y $4p=8$, así que $p=2$.

\textbf{Paso 4.} $V(3,-1)$, $F(3,1)$ y directriz $y=-3$.
\end{solution}
\end{example}

\section{Aplicaciones}

\begin{example}{Reflector parabólico}{reflector}
Un reflector tiene $12$ cm de abertura y $3$ cm de profundidad. Coloca el
vértice en el origen y encuentra la posición del foco.
\begin{solution}
\textbf{Paso 1.} Usamos $x^2=4py$. El borde contiene $(6,3)$.

\textbf{Paso 2.} Sustituimos: $36=4p(3)$, por lo que $p=3$.

\textbf{Paso 3.} El foco está a $3$ cm del vértice: $F(0,3)$.
\end{solution}
\end{example}

\section{Ejercicios y problemas}

\begin{exercise}{Parábolas}{ej-parabola}
\begin{enumerate}
  \item Obtén foco, directriz y apertura de $x^2=-16y$.
  \item Analiza $(x+2)^2=12(y-1)$.
  \item Analiza $(y-4)^2=-20(x+3)$.
  \item Obtén la ecuación con vértice $(0,0)$ y foco $(-5,0)$.
  \item Obtén la ecuación con vértice $(3,-2)$ y directriz $y=1$.
  \item Convierte $y^2+6y-4x+1=0$ a forma ordinaria.
  \item Calcula los extremos del lado recto de $(x-1)^2=8(y+2)$.
\end{enumerate}
\end{exercise}

\begin{problem}{Aplicaciones}{prob-parabola}
\begin{enumerate}
  \item Un arco parabólico mide $10$ m de ancho y $5$ m de altura. Coloca el
  vértice en la parte superior y obtén una ecuación para el arco.
  \item Una antena tiene $2$ m de diámetro y $0.25$ m de profundidad.
  Determina a qué distancia del vértice debe colocarse el receptor.
  \item Encuentra los puntos de intersección de $y=x-1$ y $x^2=4y$.
\end{enumerate}
\end{problem}

\section*{Resumen}\addcontentsline{toc}{section}{Resumen}
\begin{properties}
\[
(y-k)^2=4p(x-h),\qquad (x-h)^2=4p(y-k).
\]
Foco a distancia $|p|$ del vértice; directriz a la misma distancia en sentido
opuesto; longitud del lado recto $|4p|$.
\end{properties}
